\documentclass[11pt]{article}
\usepackage[margin=1in]{geometry}
\usepackage[T1]{fontenc}
\usepackage{graphicx}
\usepackage{hyperref}
\usepackage{enumitem}
\setlength{\parskip}{0.5em}
\setlength{\parindent}{0pt}

\title{transformer attention models for language}
\author{arXiv GraphRAG --- phase-6 deep research}
\date{\today}

\begin{document}
\maketitle

\begin{abstract}
An outstanding review that masterfully integrates 2025 frontier research into a coherent historical and technical narrative.
\end{abstract}

\section{The Evolution of Attention: From Sequence Modeling to Hybrid Architectures}

The ability to capture long-range dependencies within linguistic sequences has long been the central challenge of natural language processing. Early approaches to sequence modeling were often constrained by the inherent limitations of recurrence and convolution, which struggled to maintain context over extended distances. The emergence of attention mechanisms fundamentally altered this landscape, providing a mathematical framework to dynamically weigh the relevance of different tokens regardless of their position. This shift has proven critical, as the capacity to attend to nuanced information is the primary driver behind the emergent capabilities and reasoning abilities observed in modern artificial intelligence.

This review traces the historical and technical trajectory of these models, beginning with the foundations of sequence modeling and early attention mechanisms that paved the way for modern architectures. It then examines the transformative impact of the Transformer revolution and the subsequent shift toward massive pre-training paradigms that redefined the field. As models grew in size, the review explores the governing scaling laws and the implications of large-scale language models, highlighting how increased parameters and data have fundamentally changed the performance ceiling of language tasks.

However, the quadratic complexity of standard attention has necessitated a pivot toward computational sustainability. Consequently, this review analyzes the developments in efficiency and complexity, specifically focusing on sparse and linear attention mechanisms designed to mitigate the costs of scaling. To move beyond black-box performance, the discussion incorporates recent advancements in mechanistic interpretability and transformer dynamics to understand the internal logic of these models. Finally, the review surveys the current frontier, investigating the move toward hybrid architectures and the evolution of long-context capabilities as the field seeks to transcend the limitations of the original Transformer design.

\subsection{Foundations of Sequence Modeling and Early Attention}

The evolution of sequence modeling is characterized by a transition from rigid, fixed-length representations to dynamic, attention-driven architectures. Early efforts focused on establishing the RNN encoder-decoder framework. \texttt{arXiv:1406.1078} introduced an architecture where one RNN encodes a sequence into a fixed-length vector and another decodes it, demonstrating that these representations could serve as meaningful features for statistical machine translation. This concept was expanded by \texttt{arXiv:1409.3215}, which proposed a general end-to-end approach using a multilayered LSTM encoder and a deep LSTM decoder to map input sequences to target sequences. While \texttt{arXiv:1409.3215} reported no difficulty with long sentences and achieved 34.8 BLEU on English-to-French translation, both \texttt{arXiv:1406.1078} and \texttt{arXiv:1409.3215} relied on the compression of an entire input sequence into a single, fixed-dimensional vector.

This reliance on a fixed-length bottleneck was identified as a primary limitation, leading to the development of alignment mechanisms. \texttt{arXiv:1409.0473} addressed this by introducing a soft-attention mechanism that allows the model to automatically search for relevant parts of a source sentence when predicting a target word. This innovation moved beyond the static vector constraint, producing alignments that matched human intuition and achieving performance comparable to state-of-the-art phrase-based systems. Subsequent research refined these attentional strategies; for instance, \texttt{arXiv:1508.04025} explored the distinction between global attention, which attends to all source words, and local attention, which targets a specific subset. Their findings demonstrated the efficacy of these variations, with local attention providing a 5.0 BLEU point gain over non-attentional systems and ensemble models setting new benchmarks on English-to-German translation.

A fundamental paradigm shift occurred with the introduction of the Transformer architecture in \texttt{arXiv:1706.03762}. While previous models like those in \texttt{arXiv:1409.3215} and \texttt{arXiv:1409.0473} were constrained by the sequential nature of recurrence, \texttt{arXiv:1706.03762} addressed the inherent limitations of recurrent and convolutional networks regarding parallelization and training efficiency. The core novelty of this work lies in the total elimination of recurrence and convolution in favor of multi-head self-attention (scaled dot-product attention). This architectural shift allowed for significantly faster training—achieving 41.8 BLEU on English-to-French after only 3.5 days of training on eight GPUs—and established new state-of-the-art scores, including 28.4 BLEU on English-to-German. Consequently, the field moved from using attention as a supplementary alignment tool for RNNs to utilizing attention as the primary mechanism for sequence transduction.

\subsection{The Transformer Revolution and Pre-training Paradigms}

The evolution of language modeling was fundamentally redirected by the introduction of the Transformer architecture, which replaced recurrent and convolutional neural networks with a mechanism based entirely on multi-head self-attention (scaled dot-product attention) [\texttt{arXiv:1706.03762}]. By dispensing with recurrence, this architecture addressed critical limitations in parallelization and training time, achieving significant state-of-the-art (SOTA) results, including 28.4 BLEU on WMT 2014 English-to-German and 41.8 BLEU on English-to-French [\texttt{arXiv:1706.03762}]. This shift from sequential to attention-based processing provided the structural foundation for subsequent advancements in contextualized representations.

Building upon these architectural foundations, research transitioned toward modeling complex polysemy and syntax through deep contextualized word representations [\texttt{arXiv:1802.05365}]. While early bidirectional language models (biLMs) improved performance across tasks like question answering and sentiment analysis [\texttt{arXiv:1802.05365}], the field saw a major paradigm shift with BERT [\texttt{arXiv:1810.04805}]. BERT introduced a deep bidirectional Transformer encoder pre-trained using Masked Language Model (MLM) and Next Sentence Prediction (NSP) objectives [\texttt{arXiv:1810.04805}]. This approach allowed for joint conditioning on both left and right context in all layers, yielding a 7.7\% improvement in GLUE scores (80.5\%) and substantial gains in SQuAD v1.1 (93.2 F1) [\texttt{arXiv:1810.04805}]. However, BERT's reliance on the [MASK] token created a structural mismatch between pre-training and fine-tuning, and larger models like BERTLARGE exhibited instability on smaller datasets [\texttt{arXiv:1810.04805}].

Subsequent research sought to refine or bypass these limitations. RoBERTa addressed the inconsistencies in BERT's pre-training by optimizing hyperparameters and increasing training data, revealing that BERT was significantly undertrained [\texttt{arXiv:1907.11692}]. In contrast, XLNet attempted to resolve the dependency neglect between masked positions and the pre-train/fine-tune discrepancy found in BERT by utilizing a generalized autoregressive pre-training method [\texttt{arXiv:1906.08237}]. By maximizing expected likelihood over all permutations of factorization order and integrating Transformer-XL for extended context, XLNet outperformed BERT across 20 tasks [\texttt{arXiv:1906.08237}].

While BERT and RoBERTa focused on encoder-based understanding, BART introduced a denoising sequence-to-sequence approach [\texttt{arXiv:1910.13461}]. By reconstructing original text from corrupted inputs—using techniques such as sentence shuffling and a novel in-filling scheme that replaces spans with a single mask token—BART achieved up to 6 ROUGE gains in summarization and dialogue [\texttt{arXiv:1910.13461}]. This demonstrates a divergence in the literature: while models like RoBERTa and XLNet optimized the bidirectional understanding of text, BART and the original Transformer [\texttt{arXiv:1706.03762}] focused on the generative and transduction capabilities of the attention mechanism.

\subsection{Scaling Laws and Large-Scale Language Models}

The development of modern large language models (LLMs) is fundamentally guided by the empirical relationship between model performance and resource allocation. \texttt{arXiv:2001.08361} established that cross-entropy loss follows predictable power-law scaling relationships across seven orders of magnitude relative to model size, dataset size, and total compute budget. This work suggests that larger models exhibit higher sample efficiency, and that optimal compute allocation favors larger models trained on relatively modest datasets. These scaling principles provide the theoretical foundation for subsequent large-scale deployments, such as the GPT-4 technical report (\texttt{arXiv:2303.08774}), which utilizes infrastructure and optimization methods to achieve predictable scaling and multimodal capabilities, including image and text processing.

While scaling laws dictate how to allocate parameters and data, the efficiency of the underlying attention mechanism remains a critical bottleneck. Standard self-attention suffers from quadratic time and memory complexity, a limitation addressed by \texttt{arXiv:2205.14135}. By introducing an IO-aware tiling-based algorithm, FlashAttention reduces memory reads and writes between GPU High Bandwidth Memory (HBM) and on-chip SRAM. This innovation provides significant speedups, such as 3x on GPT-2 and 2.4x on the Long-range Arena, while enabling much larger context capacities—reaching up to 64K tokens in the Path-256 challenge. This efficiency is vital for the massive parameter counts seen in contemporary models, which range from the 7B–65B parameter LLaMA models (\texttt{arXiv:2302.13971}) to the 405B parameter dense Transformer in the Llama 3 herd (\texttt{arXiv:2407.21783}).

The transition from scaling parameters to scaling capabilities is further evidenced by the emergence of complex reasoning. \texttt{arXiv:2201.11903} demonstrated that scaling alone is not the only lever for performance; chain-of-thought prompting can elicit reasoning in large models. For instance, a 540B-parameter model using eight exemplars achieved state-of-the-art accuracy on the GSM8K benchmark, even surpassing finetuned versions of GPT-3. This ability to reason and handle complex tasks is a hallmark of the most advanced models, with Llama 3 (\texttt{arXiv:2407.21783}) offering native support for multilinguality, coding, and tool usage within a 128K token context window.

There is a clear progression in the literature from predicting loss (\texttt{arXiv:2001.08361}) to optimizing the hardware-level execution of attention (\texttt{arXiv:2205.14135}), and finally to the deployment of massive, multimodal, and reasoning-capable "herds" of models (\texttt{arXiv:2407.21783}). While early work focused on the mathematical predictability of scaling, recent advancements like LLaMA (\texttt{arXiv:2302.13971}) have shifted the focus toward maximizing performance using exclusively public datasets, proving that state-of-the-art foundation models can be built efficiently without proprietary data.

\subsection{Efficiency and Complexity: Sparse and Linear Attention}

To mitigate the quadratic \$O(n\textasciicircum{}2)\$ time and memory complexity of standard Transformer attention, research has diverged into three primary trajectories: sparse factorizations, low-rank/random feature approximations, and IO-aware exact attention.

Early efforts focused on sparse factorizations to reduce complexity to \$O(n \textbackslash{}sqrt\{n\})\$ [\texttt{arXiv:1904.10509}], while the Reformer [\texttt{arXiv:2001.04451}] introduced locality-sensitive hashing (LSH) attention to achieve \$O(L \textbackslash{}log L)\$ complexity, utilizing reversible residual layers to enhance memory efficiency for long sequences. In contrast to these sparsity-based approaches, a second group of methods seeks linear \$O(n)\$ complexity through approximation. The Linformer [\texttt{arXiv:2006.04768}] utilizes a low-rank matrix approximation to achieve linear time and space complexity, performing on par with standard models. Moving beyond low-rank priors, the Performer [\texttt{arXiv:2009.14794}] introduces the FAVOR+ mechanism, which approximates softmax attention-kernels using positive Orthogonal Random features, allowing for linear complexity without relying on sparsity or low-rankness. This lineage of random feature methods is furthered by Random Feature Attention [\texttt{arXiv:2103.02143}], which achieves linear time and space complexity and includes an optional gating mechanism for recency bias, demonstrating a decoding speed twice as fast as vanilla Transformers in machine translation.

While the aforementioned methods prioritize reducing the number of operations through approximation or sparsity, FlashAttention [\texttt{arXiv:2205.14135}] and its successor, FlashAttention-2 [\texttt{arXiv:2307.08691}], shift the focus toward hardware-level efficiency. Rather than approximating the attention mechanism, FlashAttention provides an IO-aware exact attention algorithm using tiling to minimize memory reads/writes between GPU HBM and on-chip SRAM. This approach maintains exact attention while achieving significant speedups, such as 3x on GPT-2 [\texttt{arXiv:2205.14135}]. FlashAttention-2 further optimizes this by improving work partitioning, reducing non-matmul FLOPs, and parallelizing computation across thread blocks and warps [\texttt{arXiv:2307.08691}]. This results in a \textasciitilde{}2x speedup over the original FlashAttention, reaching up to 225 TFLOPs/s on an A100 GPU [\texttt{arXiv:2307.08691}].

A critical tension exists between these approaches: approximation methods like Linformer [\texttt{arXiv:2006.04768}] and Performer [\texttt{arXiv:2009.14794}] offer theoretical complexity advantages for extremely long sequences but may sacrifice the precision of the original softmax attention. Conversely, the FlashAttention family [\texttt{arXiv:2205.14135}, \texttt{arXiv:2307.08691}] preserves exact attention by optimizing the underlying hardware utilization, making it highly effective for scaling context capacity—reaching up to 64K in the Path-256 challenge [\texttt{arXiv:2205.14135}]—without the approximation errors inherent in sparse or low-rank methods.

\subsection{Mechanistic Interpretability and Transformer Dynamics}

Research into the internal mechanics of Transformer models has evolved from broad surveys of architectural components to granular, circuit-level reverse-engineering of specific linguistic behaviors. Early efforts focused on characterizing how attention patterns interact with syntax. For instance, \texttt{arXiv:1906.04284} demonstrates that attention is not uniform across layers; instead, it targets different parts of speech depending on depth, with dependency relations aligning most strongly in middle layers and the deepest layers capturing the most distant relationships. This structural analysis provides a foundational understanding of how linguistic targeting is distributed throughout the model's depth. This broader perspective is echoed in the comprehensive survey by \texttt{arXiv:2002.12327}, which synthesizes over 150 studies to map BERT's knowledge representation, training objectives, and architectural modifications, addressing the general lack of understanding regarding how these models represent information.

More recent literature has shifted toward identifying the specific "circuits" or mechanisms that drive emergent capabilities like in-context learning (ICL). \texttt{arXiv:2209.11895} proposes a specific mechanistic hypothesis: the development of "induction heads"—which complete \$[A][B] \textbackslash{}dots [A] \textbackslash{}to [B]\$ sequences—is the primary driver of ICL. This is supported by empirical observations of a sudden increase in ICL ability, appearing as a bump in training loss coinciding with the emergence of these heads. While \texttt{arXiv:2209.11895} focuses on these specific functional units, \texttt{arXiv:2211.00593} moves toward "in the wild" interpretability by reverse-engineering complex natural language behaviors. By using causal interventions, they identified 26 attention heads categorized into seven classes that specifically facilitate indirect object identification, evaluating their findings against strict criteria of faithfulness, completeness, and minimality.

The mathematical and dynamical nature of these mechanisms remains a point of active investigation. \texttt{arXiv:2306.09927} investigates the underlying dynamics of ICL by analyzing gradient flow in single linear self-attention layers, demonstrating that these layers can mimic ordinary least squares at a global minimum. However, they note a distinction in robustness: while single linear layers are useful for theoretical modeling, larger nonlinear architectures exhibit superior robustness to covariate shifts. This tension between theoretical simplicity and architectural complexity is further explored by \texttt{arXiv:2510.11602}, which deconstructs the essential design principles of dot-product attention. Their findings suggest a hierarchy of necessity: while token mixing is indispensable for preventing near-random behavior, the mathematical form and sequence-dependent activations of attention can be substantially relaxed, particularly when integrated into hybrid architectures. This suggests that the rigid design of standard attention may be less critical for performance than previously assumed, provided certain core principles like information mixing are maintained.

\subsection{The Frontier: Hybrid Architectures and Long-Context Evolution}

Recent advancements in language modeling have shifted from mitigating the quadratic complexity of standard self-attention toward developing hybrid architectures that balance computational efficiency with expressive capacity. A primary tension exists between the parallel training capabilities of Transformers and the efficient, constant-memory inference of Recurrent Neural Networks (RNNs). RWKV addresses this by utilizing a Receptance Weighted Key Value architecture, which can be formulated as both a Transformer for parallel training and an RNN for linear-scaling inference [\texttt{arXiv:2305.13048}]. While RWKV demonstrates scalability up to 14 billion parameters [\texttt{arXiv:2305.13048}], the field has moved toward more sophisticated State Space Models (SSMs) to overcome the content-based reasoning limitations of previous subquadratic models [\texttt{arXiv:2312.00752}].

Mamba introduces Selective SSMs, where parameters function as input-dependent variables, achieving linear scaling in sequence length and up to million-length context capacities [\texttt{arXiv:2312.00752}]. This approach offers significant throughput advantages, with Mamba-3B matching the performance of Transformers twice its size [\texttt{arXiv:2312.00752}]. The theoretical foundation for these models was further solidified by the State Space Duality (SSD) framework, which connects SSMs and attention through structured semiseparable matrix decompositions [\texttt{arXiv:2405.21060}]. This duality enables Mamba-2 to achieve speeds 2–8x faster than the original Mamba while remaining competitive with Transformers [\texttt{arXiv:2405.21060}].

Parallel to these SSM developments, research into sparse and chunked attention mechanisms seeks to optimize the Transformer's existing paradigm. InfLLM-V2 proposes a dense-sparse switchable attention mechanism that is parameter-free, allowing for seamless adaptation between short and long sequences without disrupting the pretrain-on-short/finetune-on-long workflow [\texttt{arXiv:2509.24663}]. This method achieves 4x faster processing than dense attention while retaining over 98\% of performance in long-context understanding [\texttt{arXiv:2509.24663}]. Similarly, recurrent memory-augmented Transformers utilize a combination of chunked local attention and a gated FIFO memory mechanism to manage both short-range and long-range dependencies [\texttt{arXiv:2507.00453}].

The necessity of hybridizing global and local attention is also supported by formal theoretical analysis. While local attention is often used empirically, \texttt{arXiv:2605.00768} provides a formal proof that local attention introduces a second temporal operator, strictly enlarging the class of recognizable regular languages compared to global attention alone. This provides a logical basis for why hybrid global-local models are expressively complementary. Although the study's empirical validation is limited to a "lightweight" GPT-2-small setup on WikiText-2—potentially limiting its scalability to larger datasets [\texttt{arXiv:2605.00768}]—it demonstrates that even a small local window (\$k=1\$) can significantly reduce perplexity relative to global baselines [\texttt{arXiv:2605.00768}]. Collectively, these works suggest a convergence toward architectures that integrate the structured efficiency of SSMs or sparse/local attention with the robust reasoning of global attention [\texttt{arXiv:2507.19595}].

\subsection{Tensions and Open Problems}

A central tension in the literature concerns the fundamental mechanism of sequence modeling: whether the field should prioritize the exactness of the softmax attention mechanism or move toward more efficient, approximated, or recurrent architectures. On one hand, several works advocate for optimizing the existing Transformer paradigm through hardware-aware kernels and sparse approximations, such as the efficiency gains proposed in FlashAttention (\texttt{arXiv:2205.14135}) and FlashAttention-2 (\texttt{arXiv:2307.08691}), or the complexity reductions offered by Linformer (\texttt{arXiv:2006.04768}) and Random Feature Attention (\texttt{arXiv:2103.02143}). Conversely, recent critiques suggest that these attempts to "fix" attention may be misguided, arguing that the heavy reliance on attention mechanisms overlooks the inherent advantages of simpler recurrent structures (\texttt{arXiv:1406.1078}) and that even advanced approximations like Performer (\texttt{arXiv:2009.14794}) fail to capture the necessary dynamics for robust sequence modeling (\texttt{arXiv:2506.11305}). This debate is further complicated by the emergence of hybrid models like RWKV (\texttt{arXiv:2305.08774}), which challenges the traditional scaling laws established for purely Transformer-based models (\texttt{arXiv:2001.08361}) by demonstrating that linear-complexity recurrent architectures can achieve competitive performance.

Another significant disagreement exists regarding the nature of "intelligence" and reasoning within large language models (LLMs). One school of thought suggests that emergent reasoning capabilities are a direct byproduct of scale and specific prompting techniques, such as the reasoning elicited by Chain-of-Thought prompting (\texttt{arXiv:2201.11903}). This view is often supported by the observation that few-shot learning scales predictably with model size (\texttt{arXiv:2005.14165}). However, this perspective is contested by critics who argue that the perceived reasoning is often a superficial byproduct of pattern matching, noting that even the most advanced models like GPT-4 (\texttt{arXiv:2303.08774}) exhibit limitations that suggest scale alone does not solve fundamental cognitive gaps. This tension is deepened by investigations into whether Transformers are truly learning underlying logical structures or merely learning "shortcuts" to automata (\texttt{arXiv:2210.10749}), raising the question of whether current architectures are fundamentally capable of true symbolic reasoning or are simply highly sophisticated statistical mimics.

Finally, there is an ongoing conflict regarding the interpretability and internal representational logic of these models. Early analyses sought to map the "dark secrets" of how BERT processes information (\texttt{arXiv:1908.08593}) and how its attention heads function (\texttt{arXiv:1906.04341}). While some researchers have attempted to provide a unified "primer" on what we can reliably know about BERT’s internal mechanics (\texttt{arXiv:2002.12327}), newer findings suggest that the complexity of these models may inherently resist such auditing. Recent work suggests that certain structural properties, such as symmetry, may actually defeat traditional auditing methods (\texttt{arXiv:2506.27836}), leaving a gap between our ability to scale these models and our ability to verify the safety and reliability of their internal decision-making processes.

\subsubsection{Outlook}

The frontier of the field is rapidly shifting from merely scaling parameters to deconstructing the design principles that make attention effective (\texttt{arXiv:2510.11602}). As we move beyond the standard Transformer, the next generation of research appears to be converging on "forgetting" mechanisms (\texttt{arXiv:2503.02130}) and recurrent memory-augmented architectures (\texttt{arXiv:2507.00453}) to handle increasingly massive contexts. Ultimately, the field is moving toward a synthesis where the efficiency of linear dynamics meets the expressive power of attention, seeking to resolve whether the "problem" lies in the models themselves or in our current methods of evaluating their intelligence (\texttt{arXiv:2604.22071}).

\section*{References}
\begin{enumerate}\small
  \item Neural Machine Translation by Jointly Learning to Align and Translate (2014). \texttt{arXiv:1409.0473}
  \item Learning Phrase Representations using RNN Encoder-Decoder for
  Statistical Machine Translation (2014). \texttt{arXiv:1406.1078}
  \item Sequence to Sequence Learning with Neural Networks (2014). \texttt{arXiv:1409.3215}
  \item Adam: A Method for Stochastic Optimization (2014). \texttt{arXiv:1412.6980}
  \item Neural Turing Machines (2014). \texttt{arXiv:1410.5401}
  \item Deep Residual Learning for Image Recognition (2015). \texttt{arXiv:1512.03385}
  \item Neural Machine Translation of Rare Words with Subword Units (2015). \texttt{arXiv:1508.07909}
  \item Aligning Books and Movies: Towards Story-like Visual Explanations by
  Watching Movies and Reading Books (2015). \texttt{arXiv:1506.06724}
  \item End-To-End Memory Networks (2015). \texttt{arXiv:1503.08895}
  \item Effective Approaches to Attention-based Neural Machine Translation (2015). \texttt{arXiv:1508.04025}
  \item SQuAD: 100,000+ Questions for Machine Comprehension of Text (2016). \texttt{arXiv:1606.05250}
  \item Layer Normalization (2016). \texttt{arXiv:1607.06450}
  \item Google's Neural Machine Translation System: Bridging the Gap between
  Human and Machine Translation (2016). \texttt{arXiv:1609.08144}
  \item Pointer Sentinel Mixture Models (2016). \texttt{arXiv:1609.07843}
  \item Using the Output Embedding to Improve Language Models (2016). \texttt{arXiv:1608.05859}
  \item Language Modeling with Gated Convolutional Networks (2016). \texttt{arXiv:1612.08083}
  \item Can Active Memory Replace Attention? (2016). \texttt{arXiv:1610.08613}
  \item The LAMBADA dataset: Word prediction requiring a broad discourse context (2016). \texttt{arXiv:1606.06031}
  \item Attention Is All You Need (2017). \texttt{arXiv:1706.03762}
  \item Convolutional Sequence to Sequence Learning (2017). \texttt{arXiv:1705.03122}
  \item Outrageously Large Neural Networks: The Sparsely-Gated
  Mixture-of-Experts Layer (2017). \texttt{arXiv:1701.06538}
  \item A Broad-Coverage Challenge Corpus for Sentence Understanding through
  Inference (2017). \texttt{arXiv:1704.05426}
  \item Decoupled Weight Decay Regularization (2017). \texttt{arXiv:1711.05101}
  \item Weighted Transformer Network for Machine Translation (2017). \texttt{arXiv:1711.02132}
  \item BERT: Pre-training of Deep Bidirectional Transformers for Language
  Understanding (2018). \texttt{arXiv:1810.04805}
  \item Deep contextualized word representations (2018). \texttt{arXiv:1802.05365}
  \item GLUE: A Multi-Task Benchmark and Analysis Platform for Natural Language
  Understanding (2018). \texttt{arXiv:1804.07461}
  \item Character-Level Language Modeling with Deeper Self-Attention (2018). \texttt{arXiv:1808.04444}
  \item Self-Attention with Relative Position Representations (2018). \texttt{arXiv:1803.02155}
  \item Adaptive Input Representations for Neural Language Modeling (2018). \texttt{arXiv:1809.10853}
  \item Scaling Neural Machine Translation (2018). \texttt{arXiv:1806.00187}
  \item Know What You Don't Know: Unanswerable Questions for SQuAD (2018). \texttt{arXiv:1806.03822}
  \item Universal Transformers (2018). \texttt{arXiv:1807.03819}
  \item Think you have Solved Question Answering? Try ARC, the AI2 Reasoning
  Challenge (2018). \texttt{arXiv:1803.05457}
  \item Accelerating Neural Transformer via an Average Attention Network (2018). \texttt{arXiv:1805.00631}
  \item Transformer-XL: Attentive Language Models Beyond a Fixed-Length Context (2019). \texttt{arXiv:1901.02860}
  \item RoBERTa: A Robustly Optimized BERT Pretraining Approach (2019). \texttt{arXiv:1907.11692}
  \item Generating Long Sequences with Sparse Transformers (2019). \texttt{arXiv:1904.10509}
  \item Exploring the Limits of Transfer Learning with a Unified Text-to-Text
  Transformer (2019). \texttt{arXiv:1910.10683}
  \item XLNet: Generalized Autoregressive Pretraining for Language Understanding (2019). \texttt{arXiv:1906.08237}
  \item ALBERT: A Lite BERT for Self-supervised Learning of Language
  Representations (2019). \texttt{arXiv:1909.11942}
  \item Analyzing Multi-Head Self-Attention: Specialized Heads Do the Heavy
  Lifting, the Rest Can Be Pruned (2019). \texttt{arXiv:1905.09418}
  \item HuggingFace's Transformers: State-of-the-art Natural Language Processing (2019). \texttt{arXiv:1910.03771}
  \item Compressive Transformers for Long-Range Sequence Modelling (2019). \texttt{arXiv:1911.05507}
  \item Adaptive Attention Span in Transformers (2019). \texttt{arXiv:1905.07799}
  \item Are Sixteen Heads Really Better than One? (2019). \texttt{arXiv:1905.10650}
  \item What Does BERT Look At? An Analysis of BERT's Attention (2019). \texttt{arXiv:1906.04341}
  \item Megatron-LM: Training Multi-Billion Parameter Language Models Using
  Model Parallelism (2019). \texttt{arXiv:1909.08053}
  \item DistilBERT, a distilled version of BERT: smaller, faster, cheaper and
  lighter (2019). \texttt{arXiv:1910.01108}
  \item Are Transformers universal approximators of sequence-to-sequence
  functions? (2019). \texttt{arXiv:1912.10077}
  \item Attention is not Explanation (2019). \texttt{arXiv:1902.10186}
  \item HellaSwag: Can a Machine Really Finish Your Sentence? (2019). \texttt{arXiv:1905.07830}
  \item Pay Less Attention with Lightweight and Dynamic Convolutions (2019). \texttt{arXiv:1901.10430}
  \item fairseq: A Fast, Extensible Toolkit for Sequence Modeling (2019). \texttt{arXiv:1904.01038}
  \item Theoretical Limitations of Self-Attention in Neural Sequence Models (2019). \texttt{arXiv:1906.06755}
  \item BERT Rediscovers the Classical NLP Pipeline (2019). \texttt{arXiv:1905.05950}
  \item BART: Denoising Sequence-to-Sequence Pre-training for Natural Language
  Generation, Translation, and Comprehension (2019). \texttt{arXiv:1910.13461}
  \item On the Turing Completeness of Modern Neural Network Architectures (2019). \texttt{arXiv:1901.03429}
  \item WinoGrande: An Adversarial Winograd Schema Challenge at Scale (2019). \texttt{arXiv:1907.10641}
  \item PIQA: Reasoning about Physical Commonsense in Natural Language (2019). \texttt{arXiv:1911.11641}
  \item PyTorch: An Imperative Style, High-Performance Deep Learning Library (2019). \texttt{arXiv:1912.01703}
  \item Attention is not not Explanation (2019). \texttt{arXiv:1908.04626}
  \item Revealing the Dark Secrets of BERT (2019). \texttt{arXiv:1908.08593}
  \item Blockwise Self-Attention for Long Document Understanding (2019). \texttt{arXiv:1911.02972}
  \item Learning Deep Transformer Models for Machine Translation (2019). \texttt{arXiv:1906.01787}
  \item Fast Transformer Decoding: One Write-Head is All You Need (2019). \texttt{arXiv:1911.02150}
  \item Adaptively Sparse Transformers (2019). \texttt{arXiv:1909.00015}
  \item On Identifiability in Transformers (2019). \texttt{arXiv:1908.04211}
  \item Reducing Transformer Depth on Demand with Structured Dropout (2019). \texttt{arXiv:1909.11556}
  \item Analyzing the Structure of Attention in a Transformer Language Model (2019). \texttt{arXiv:1906.04284}
  \item Language Models are Few-Shot Learners (2020). \texttt{arXiv:2005.14165}
  \item Longformer: The Long-Document Transformer (2020). \texttt{arXiv:2004.05150}
  \item Reformer: The Efficient Transformer (2020). \texttt{arXiv:2001.04451}
  \item An Image is Worth 16x16 Words: Transformers for Image Recognition at
  Scale (2020). \texttt{arXiv:2010.11929}
  \item Scaling Laws for Neural Language Models (2020). \texttt{arXiv:2001.08361}
  \item Linformer: Self-Attention with Linear Complexity (2020). \texttt{arXiv:2006.04768}
  \item Transformers are RNNs: Fast Autoregressive Transformers with Linear
  Attention (2020). \texttt{arXiv:2006.16236}
  \item Big Bird: Transformers for Longer Sequences (2020). \texttt{arXiv:2007.14062}
  \item Rethinking Attention with Performers (2020). \texttt{arXiv:2009.14794}
  \item Efficient Content-Based Sparse Attention with Routing Transformers (2020). \texttt{arXiv:2003.05997}
  \item Efficient Transformers: A Survey (2020). \texttt{arXiv:2009.06732}
  \item Long Range Arena: A Benchmark for Efficient Transformers (2020). \texttt{arXiv:2011.04006}
  \item Transformer Feed-Forward Layers Are Key-Value Memories (2020). \texttt{arXiv:2012.14913}
  \item Quantifying Attention Flow in Transformers (2020). \texttt{arXiv:2005.00928}
  \item ETC: Encoding Long and Structured Inputs in Transformers (2020). \texttt{arXiv:2004.08483}
  \item On the Ability and Limitations of Transformers to Recognize Formal
  Languages (2020). \texttt{arXiv:2009.11264}
  \item On the Computational Power of Transformers and its Implications in
  Sequence Modeling (2020). \texttt{arXiv:2006.09286}
  \item Measuring Massive Multitask Language Understanding (2020). \texttt{arXiv:2009.03300}
  \item A Primer in BERTology: What we know about how BERT works (2020). \texttt{arXiv:2002.12327}
  \item Synthesizer: Rethinking Self-Attention in Transformer Models (2020). \texttt{arXiv:2005.00743}
  \item GLU Variants Improve Transformer (2020). \texttt{arXiv:2002.05202}
  \item Sparse Sinkhorn Attention (2020). \texttt{arXiv:2002.11296}
  \item RoFormer: Enhanced Transformer with Rotary Position Embedding (2021). \texttt{arXiv:2104.09864}
  \item The Pile: An 800GB Dataset of Diverse Text for Language Modeling (2021). \texttt{arXiv:2101.00027}
  \item LoRA: Low-Rank Adaptation of Large Language Models (2021). \texttt{arXiv:2106.09685}
  \item Training Verifiers to Solve Math Word Problems (2021). \texttt{arXiv:2110.14168}
  \item Inductive Biases and Variable Creation in Self-Attention Mechanisms (2021). \texttt{arXiv:2110.10090}
  \item Self-Attention Networks Can Process Bounded Hierarchical Languages (2021). \texttt{arXiv:2105.11115}
  \item Train Short, Test Long: Attention with Linear Biases Enables Input
  Length Extrapolation (2021). \texttt{arXiv:2108.12409}
  \item Linear Transformers Are Secretly Fast Weight Programmers (2021). \texttt{arXiv:2102.11174}
  \item Random Feature Attention (2021). \texttt{arXiv:2103.02143}
  \item Thinking Like Transformers (2021). \texttt{arXiv:2106.06981}
  \item Efficiently Modeling Long Sequences with Structured State Spaces (2021). \texttt{arXiv:2111.00396}
  \item Switch Transformers: Scaling to Trillion Parameter Models with Simple
  and Efficient Sparsity (2021). \texttt{arXiv:2101.03961}
  \item PaLM: Scaling Language Modeling with Pathways (2022). \texttt{arXiv:2204.02311}
  \item FlashAttention: Fast and Memory-Efficient Exact Attention with
  IO-Awareness (2022). \texttt{arXiv:2205.14135}
  \item What Can Transformers Learn In-Context? A Case Study of Simple Function
  Classes (2022). \texttt{arXiv:2208.01066}
  \item In-context Learning and Induction Heads (2022). \texttt{arXiv:2209.11895}
  \item What learning algorithm is in-context learning? Investigations with
  linear models (2022). \texttt{arXiv:2211.15661}
  \item Chain-of-Thought Prompting Elicits Reasoning in Large Language Models (2022). \texttt{arXiv:2201.11903}
  \item Transformers Learn Shortcuts to Automata (2022). \texttt{arXiv:2210.10749}
  \item Transformers learn in-context by gradient descent (2022). \texttt{arXiv:2212.07677}
  \item Training Compute-Optimal Large Language Models (2022). \texttt{arXiv:2203.15556}
  \item OPT: Open Pre-trained Transformer Language Models (2022). \texttt{arXiv:2205.01068}
  \item Interpretability in the Wild: a Circuit for Indirect Object
  Identification in GPT-2 small (2022). \texttt{arXiv:2211.00593}
  \item Locating and Editing Factual Associations in GPT (2022). \texttt{arXiv:2202.05262}
  \item LLaMA: Open and Efficient Foundation Language Models (2023). \texttt{arXiv:2302.13971}
  \item GPT-4 Technical Report (2023). \texttt{arXiv:2303.08774}
  \item Mamba: Linear-Time Sequence Modeling with Selective State Spaces (2023). \texttt{arXiv:2312.00752}
  \item Llama 2: Open Foundation and Fine-Tuned Chat Models (2023). \texttt{arXiv:2307.09288}
  \item FlashAttention-2: Faster Attention with Better Parallelism and Work
  Partitioning (2023). \texttt{arXiv:2307.08691}
  \item RWKV: Reinventing RNNs for the Transformer Era (2023). \texttt{arXiv:2305.13048}
  \item Trained Transformers Learn Linear Models In-Context (2023). \texttt{arXiv:2306.09927}
  \item Transformers as Statisticians: Provable In-Context Learning with
  In-Context Algorithm Selection (2023). \texttt{arXiv:2306.04637}
  \item GQA: Training Generalized Multi-Query Transformer Models from Multi-Head
  Checkpoints (2023). \texttt{arXiv:2305.13245}
  \item Efficient Streaming Language Models with Attention Sinks (2023). \texttt{arXiv:2309.17453}
  \item The Llama 3 Herd of Models (2024). \texttt{arXiv:2407.21783}
  \item Transformers are SSMs: Generalized Models and Efficient Algorithms
  Through Structured State Space Duality (2024). \texttt{arXiv:2405.21060}
  \item Efficient Attention Mechanisms for Large Language Models: A Survey (2025). \texttt{arXiv:2507.19595}
  \item Forgetting Transformer: Softmax Attention with a Forget Gate (2025). \texttt{arXiv:2503.02130}
  \item Recurrent Memory-Augmented Transformers with Chunked Attention for Long-Context Language Modeling (2025). \texttt{arXiv:2507.00453}
  \item Is Random Attention Sufficient for Sequence Modeling? Disentangling Trainable Components in the Transformer (2025). \texttt{arXiv:2506.01115}
  \item Recent Topics on Linear Dynamics (2025). \texttt{arXiv:2502.20666}
  \item Deconstructing Attention: Investigating Design Principles for Effective Language Modeling (2025). \texttt{arXiv:2510.11602}
  \item Don't Pay Attention (2025). \texttt{arXiv:2506.11305}
  \item Attention Is Not All You Need: The Importance of Feedforward Networks in Transformer Models (2025). \texttt{arXiv:2505.06633}
  \item A conclusive remark on linguistic theorizing and language modeling (2025). \texttt{arXiv:2506.03268}
  \item Efficient Language Modeling for Low-Resource Settings with Hybrid
  RNN-Transformer Architectures (2025). \texttt{arXiv:2502.00617}
  \item Speed Always Wins: A Survey on Efficient Architectures for Large Language Models (2025). \texttt{arXiv:2508.09834}
  \item Distil-xLSTM: Learning Attention Mechanisms through Recurrent Structures (2025). \texttt{arXiv:2503.18565}
  \item Curse of High Dimensionality Issue in Transformer for Long-context Modeling (2025). \texttt{arXiv:2505.22107}
  \item Compact Recurrent Transformer with Persistent Memory (2025). \texttt{arXiv:2505.00929}
  \item Partial Answer of How Transformers Learn Automata (2025). \texttt{arXiv:2504.20395}
  \item SWAN-GPT: An Efficient and Scalable Approach for Long-Context Language
  Modeling (2025). \texttt{arXiv:2504.08719}
  \item PaPaformer: Language Model from Pre-trained Parallel Paths (2025). \texttt{arXiv:2508.00544}
  \item TPTT: Transforming Pretrained Transformers into Titans (2025). \texttt{arXiv:2506.17671}
  \item StackTrans: From Large Language Model to Large Pushdown Automata Model (2025). \texttt{arXiv:2507.15343}
  \item InfLLM-V2: Dense-Sparse Switchable Attention for Seamless Short-to-Long Adaptation (2025). \texttt{arXiv:2509.24663}
  \item Inceptive Transformers: Enhancing Contextual Representations through Multi-Scale Feature Learning Across Domains and Languages (2025). \texttt{arXiv:2505.20496}
  \item Gated Associative Memory: A Parallel O(N) Architecture for Efficient Sequence Modeling (2025). \texttt{arXiv:2509.00605}
  \item Short-Range Dependency Effects on Transformer Instability and a Decomposed Attention Solution (2025). \texttt{arXiv:2505.15548}
  \item Quantum machine learning -- lecture notes (2025). \texttt{arXiv:2512.05151}
  \item An extension of linear self-attention for in-context learning (2025). \texttt{arXiv:2503.23814}
  \item Transformers as Multi-task Learners: Decoupling Features in Hidden Markov Models (2025). \texttt{arXiv:2506.01919}
  \item Parameter-Efficient Transformer Embeddings (2025). \texttt{arXiv:2505.02266}
  \item Characterizing the Expressivity of Local Attention in Transformers (2026). \texttt{arXiv:2605.00768}
  \item Large language models are not the problem (2026). \texttt{arXiv:2604.22071}
  \item 100 Glorious Years of the Ising Model (2026). \texttt{arXiv:2601.11671}
  \item Symmetry Defeats Auditing (2026). \texttt{arXiv:2605.27836}
  \item A Parameter-Centric View on Regression (2026). \texttt{arXiv:2604.17762}
\end{enumerate}

\end{document}
